\documentclass[tikz,border=8pt]{standalone}
\usepackage{tikz}
\usepackage{amsmath}
\usepackage{pgfplots}
\pgfplotsset{compat=1.18}
\usepackage{ctex}
\usetikzlibrary{calc, positioning, arrows.meta, shapes, fit, backgrounds}

\begin{document}
\begin{tikzpicture}[
  >=Stealth,
  node dist/.style={node distance=2.5cm},
  bn node/.style={circle, draw=blue!70!black, fill=blue!10,
                  minimum size=14mm, font=\normalsize\bfseries, thick},
  cpt box/.style={rectangle, draw=gray!60, fill=white, thin,
                  rounded corners=2pt, font=\scriptsize, inner sep=4pt, align=left}
]

% 节点布局
%   A
%  / \
% B   C
%  \ /
%   D

\node[bn node] (A) at (4.5, 7.0) {$A$};
\node[bn node] (B) at (2.0, 4.5) {$B$};
\node[bn node] (C) at (7.0, 4.5) {$C$};
\node[bn node] (D) at (4.5, 2.0) {$D$};

% 有向边
\draw[->, thick, blue!60!black] (A) -- (B);
\draw[->, thick, blue!60!black] (A) -- (C);
\draw[->, thick, blue!60!black] (B) -- (D);
\draw[->, thick, blue!60!black] (C) -- (D);

% 节点旁的条件概率表（CPT）
% A 的 CPT（先验）
\node[cpt box] at (-0.5, 7.0) {
  \textbf{P(A)}\\
  P(A=0)=0.4\\
  P(A=1)=0.6
};
\draw[gray!40, thin] (-0.5+1.2, 7.0) -- (A.west);

% B 的 CPT
\node[cpt box] at (-0.5, 4.5) {
  \textbf{P(B|A)}\\
  P(B=1|A=0)=0.2\\
  P(B=1|A=1)=0.8
};
\draw[gray!40, thin] (-0.5+1.8, 4.5) -- (B.west);

% C 的 CPT
\node[cpt box] at (9.5, 4.5) {
  \textbf{P(C|A)}\\
  P(C=1|A=0)=0.3\\
  P(C=1|A=1)=0.7
};
\draw[gray!40, thin] (9.5-1.3, 4.5) -- (C.east);

% D 的 CPT
\node[cpt box] at (4.5, -0.3) {
  \textbf{P(D|B,C)}\\
  P(D=1|B=0,C=0)=0.1\quad P(D=1|B=0,C=1)=0.5\\
  P(D=1|B=1,C=0)=0.6\quad P(D=1|B=1,C=1)=0.9
};
\draw[gray!40, thin] (4.5, -0.3+0.5) -- (D.south);

% 整体标题
\node[font=\large\bfseries] at (4.5, 8.6) {贝叶斯网络（DAG）};

% 联合概率分解公式
\node[font=\normalsize,
      fill=white, draw=gray!50, thin, rounded corners=2pt, inner sep=6pt]
  at (4.5, -2.2) {
    $P(A,B,C,D)=P(A)\,P(B\mid A)\,P(C\mid A)\,P(D\mid B,C)$
  };

\end{tikzpicture}
\end{document}
